\documentclass[../ICIT-Paper-2024.tex]{subfiles}
\usepackage{hyperref} 
\usepackage{tablefootnote}
\begin{document}
\section{EXPERIMENT}\label{section4}
Before applying any method, the original dataset was divided into two parts: (1) the training dataset and (2) the testing dataset. The training dataset is used to train the model, while the testing dataset is used to evaluate the accuracy of the model. This ensures that the model not only memorizes the data it has seen but also has good generalization ability on new data. Typically, the dataset is split in an 80:20 ratio, with 80\% for training and 20\% for testing. However, this ratio may vary depending on the specific problem and can be adjusted to achieve the best results.

In this section, the results of the methods are presented and evaluated using two main performance metrics:

\textbf{Accuracy:} Accuracy is the most widely used classification performance metric. It calculates the number of correctly classified samples by the model out of the total samples in the dataset. Accuracy is calculated as explained in Equation.~\ref{eq3}

\begin{equation}\label{eq3}
    \text{Accuracy} = \frac{\text{Count of Correct Predictions}}{\text{Total Number of Predictions}}
\end{equation}

\textbf{Precision:} Precision is a metric used to evaluate the accuracy of correctly classified positive samples by the model. It is calculated by dividing the number of TP by the sum of TP and FP. In simpler terms, precision measures how many of the classified TP instances are actually positive. Precision is calculated as explained in Equation~\ref{eq4}

\begin{equation}\label{eq4}
    \text{Precision} = \frac{\text{TP}}{\text{TP} + \text{FP}}
\end{equation}

In some cases, each prediction may have more than one correct label. For instance, a wallet, which has a borrow-to-asset ratio over 40\% and more than one liquidation, can be Poor or Fair. As a result, Recall (Detection Rate) cannot be calculated. Therefore, Recall should not be used to evaluate the performance of the models in this case.



In summary, 

\subsection{Experiment 1}
\textbf{SGD, Adam: }
% \subsubsection{SGD, Adam}
% \paragraph{}
After determining the evaluation method, we proceeded to implement the models. For SGD, we initialized the $decay\ rate = 0.95$. The corresponding parameters for Adam include $\beta_1 = 0.9$, $\beta_2 = 0.999$, and $\epsilon = 1\times10^{-8}$. Both algorithms were tested with different values of learning rate and number of iterations. Specifically, the learning rates included values of $0.1$, $0.01$, and $0.001$, and the number of iterations included values of $100$, $1000$, and $10000$. After implementation, both SGD and Adam algorithms yielded the best results at a learning rate of $0.001$ and number of iterations of $10000$.

%Sau khi xác định được phương pháp đánh giá, chúng tôi tiến hành triển khai các models. Đối với SGD, chúng tôi khởi tạo than số  decay rate là 0.95. Các tham số  tương ứng đối với Adam bao gồm beta1 là 0.9, beta2 là 0.999, và epsilon là 1e-8. Cả hai thuật toán đều được thử nghiệm với các giá trị learning rate và number iterations khác nhau. Cụ thể, learning rate bao gồm  các giá trị là 0.1, 0.01, 0.001, và number iterations bao gồm  các giá trị 100, 1000, 10000. Sau khi triển khai, cả 2 thuật toán SGD, Adam cho kết quả tốt nhất tại learning rate là 0.001 và number iterations là 10000.

\textbf{Genetic Algorithm: }
% \subsubsection{Genetic Algorithm}
% \paragraph{}
For the GA algorithm, we initialized the number of generations parameter to $100$. From this parameter, we tested the model with different crossover functions and different mutation rates. Mutation rates of $0.1$ and $0.2$, and crossover functions of one-point and two-point were used. After implementation and evaluation, the two-point crossover function combined with a mutation rate of $0.1$ gives the best results.

%Đối với giải thuât GA, chúng tôi khởi tạo tham số number generations là 100. Từ tham số này, chúng tôi thử nghiệm mô hình với các hàm lai tạo khác nhau với các tỉ lệ đột biến khác nhau. Tỉ lệ đột biến 0.1, 0.2, và hàm lai tạo one point, two point được chúng tôi sử dụng. Sau khi triển khai và đánh giá thì hàm lai tạo two point kêt hợp với tỉ lệ đột biến 0.1 cho kết quả tốt nhất.
\textbf{Multilayer Perception: }
% \subsubsection{Multilayer Perception}
% \paragraph{}
For Multilayer Perceptron, we tested the model with different neural network architectures, as well as varying the number of epochs and batch sizes. Specifically, the values for epochs included 5, 10, and 20, while the values for batch size included 32 and 64. After implementation, the model with 2 hidden layers, as introduced above, along with 10 epochs and a batch size of 64, gives the best results.


The results of our study are organized in the tables below (Tables ~\ref{tab:results})


\begin{table}[h!]
\centering
\caption{The performance of the optimized models}\label{tab:results}
\begin{tabular}{lllll}
\hline
\textbf{Model}                      & \textbf{Accuracy}         & \textbf{Precision}  \\
                                    & (\%)          & (\%)          \\ \hline
SGD                                 & 76.56         & 35.11         \\
Adam                                & 95.12         & 82.50         \\
Genetic Algorithm                   & 96.94         & 85.78         \\
Multilayer Perception               & 97.50         & 91.65         \\
\hline
\\
\end{tabular}

\end{table}

\begin{figure}[h!]
    \centering
    \includegraphics[scale=0.5]{heatmap.pdf}
    \caption{Caption of the Heatmap}
    \label{fig:heatmap}
\end{figure}

\end{document}